% Generated by Listas/make_statements.py from pratica-lecture-07.md, item 2. Do not edit.
\textbf{Source.} Muñoz-Garcia \& Toro-Gonzalez (2019), \emph{Strategy and Game Theory: Practice Exercises with Answers}, Ch. 7, Exercise 1 ``A Simple Bayesian Game'', pp.~318--319 {[}PDF pp.~329--330{]}. Reproduced verbatim (statement only; the book's worked answer follows it). \textbf{Answer in the source:} yes, the book's worked answer follows the exercise.



\begin{quote}

Consider the payoff matrix in Fig. 7.1, where player 1 (in rows) chooses between \(A\) and \(B\), and player 2 (in columns) selects \(C\) or \(D\). Player 1 observes the realization of random variable \(x\), which takes value \(x = 15\) with probability 1/2 or \(x = 4\) with probability 1/2. Player 2, however, doesn't observe the realization of \(x\), but knows its probability distribution. Therefore, player 1 is perfectly informed about the payoff matrix they face, while player 2 does not know the payoff that player 1 earns in strategy profile \((A, C)\). While this payoff only affects player 1, player 2's behavior is affected by this uncertainty since he could anticipate that player 1 best responds to column \(C\) with \(A\) when \(x = 15\), but player 1 would best respond with \(B\) when \(x = 4\).



{\def\LTcaptype{none} % do not increment counter

\begin{longtable}[]{@{}lll@{}}

\toprule\noalign{}

Player 1 ~Player 2 & \emph{C} & \emph{D} \\

\midrule\noalign{}

\endhead

\bottomrule\noalign{}

\endlastfoot

\emph{A} & \(x\), 18 & 6, 7 \\

\emph{B} & 8, 0 & 8, 11 \\

\end{longtable}

}



\emph{Fig. 7.1} Simultaneous-move game under asymmetric information about \(x\)



\begin{enumerate}

\def\labelenumi{(\alph{enumi})}

\item

  \emph{Complete information}. As a benchmark, let us assume in this part of the exercise that both players are informed about the realization of \(x\). Which are the pure strategy Nash equilibria of the game when \(x = 15\)? And what are the pure strategy Nash equilibria when \(x = 4\)?

\item

  \emph{Incomplete information}. Let us now return to the game where only player 1 privately observes the realization of \(x\). Find the Bayesian Nash Equilibrium of this incomplete information game by following the next three steps:

\end{enumerate}



\begin{enumerate}

\def\labelenumi{\arabic{enumi}.}

\tightlist

\item

  Identify each player's strategy space in this incomplete information setting,

\item

  Construct the Bayesian normal form representation, and

\item

  Underline each player's best response payoffs to find the BNE of the game.

\end{enumerate}

\end{quote}

